%% toc: P3 38 38.7
\section{Proofs for the comparison of regimes}
\label{sec:p03-w03-p03-053-annex}

\begin{annexitem}{annex:p03-w03-p03-053-code-carrier}{thm:p03-w03-p03-053-code-carrier}{Same code, different carrier topologies}
\begin{proof}
For (1), the two topological structures have the same carrier and the
same interpretation of every symbol of \(L\). Satisfaction of an
\(L_{\omega_1\omega}\)-formula of \(L\) is defined by induction through
atoms of \(L\), negation, countable conjunction and disjunction, and
point quantifiers over the carrier. None of these clauses consults the
open sort, so an induction on formulas shows that both structures satisfy
the same such formulas under the same point assignments. Their
\(L\)-reduct is \(\mathcal M\) itself, so they correspond to the same
code.

For (2), let \(X\) be the carrier; it has at least two points. In the
discrete topology every singleton is open. For distinct \(x,y\) the open
set \(U=\{x\}\) contains \(x\) and not \(y\), which witnesses the
guarded existential in \(T_1\); hence \(T_1\) holds. In the indiscrete
topology the open sets are \(\varnothing\) and \(X\). Choose distinct
\(x,y\in X\). The only open set containing \(x\) is \(X\), which contains
\(y\), so the first disjunct of \(T_0\) fails for this pair; by symmetry
the second fails as well. Hence \(T_0\) is false.

Thus two topological structures over one code disagree on an
\(L_t\)-sentence, and \(L_t\)-truth is not a function of the code. The
topology of \(X_L\) is the same for both structures by (1), so it
determines neither carrier topology.
\end{proof}
\end{annexitem}

\begin{annexitem}{annex:p03-w03-p03-053-almost-sure}{thm:p03-w03-p03-053-almost-sure}{Completeness, satisfiability and infinity of the almost-sure theory}
\begin{proof}
For a sentence \(\varphi\), the events ``\(\varphi\) holds'' and
``\(\neg\varphi\) holds'' partition the finite sample space of size-\(n\)
structures, so \(\mu_n(\neg\varphi)=1-\mu_n(\varphi)\). By
\cref{thm:p03-fol-zero-one}, \(\lim_n\mu_n(\varphi)\) exists and is
\(0\) or \(1\). If it is \(1\), then \(\varphi\in T_{\mathrm{as}}\) and the
limit for \(\neg\varphi\) is \(0\); if it is \(0\), the reverse holds.
So exactly one of \(\varphi,\neg\varphi\) lies in \(T_{\mathrm{as}}\).

Let \(\varphi_1,\ldots,\varphi_k\in T_{\mathrm{as}}\). The probability
that some \(\varphi_i\) fails is at most
\(\sum_{i=1}^k\mu_n(\neg\varphi_i)\), which tends to \(0\). Hence for all
sufficiently large \(n\) the probability that all \(\varphi_i\) hold is
positive, and some structure with universe \([n]\) satisfies all of them.
Every finite subset of \(T_{\mathrm{as}}\) is therefore satisfiable, and
first-order compactness (\cref{thm:fol-compactness}) gives a model
of \(T_{\mathrm{as}}\).

For \(k\ge1\) let \(\varepsilon_k\) be
\(\exists x_1\cdots\exists x_k\bigwedge_{1\le i<j\le k}x_i\neq x_j\).
Every structure with universe \([n]\), \(n\ge k\), satisfies
\(\varepsilon_k\), so \(\mu_n(\varepsilon_k)=1\) for \(n\ge k\) and
\(\varepsilon_k\in T_{\mathrm{as}}\). A model of all \(\varepsilon_k\)
has at least \(k\) elements for every \(k\), so it is infinite.

Finally let \(\mathcal A\) have universe \([n]\). Then
\(\mathcal A\) satisfies \(\neg\varepsilon_{n+1}\), while
\[
 \lim_{m\to\infty}\mu_m(\neg\varepsilon_{n+1})=0,
\]
because \(\varepsilon_{n+1}\in T_{\mathrm{as}}\). This sentence has truth value
one in \(\mathcal A\) and limit probability zero, so
\(T_{\mathrm{as}}\) is not the theory of \(\mathcal A\).
\end{proof}
\end{annexitem}

\begin{annexitem}{annex:p03-w03-p03-053-finite-collapse}{thm:p03-w03-p03-053-finite-collapse}{Finite carriers}
\begin{proof}
A Polish space is metrizable, so distinct points have disjoint open
neighbourhoods and every singleton is closed. In a finite space the
complement of a singleton is a finite union of closed singletons, hence
closed, so every singleton is open and the topology is discrete.

In a discrete space every set is both open and closed. A nowhere dense
set \(A\) satisfies \(\varnothing=\operatorname{int}\operatorname{cl}A=A\).
A meager set is contained in a countable union of empty sets, so it is
empty, and the only comeager set is \(X\). The section
\(\varphi(\mathcal M,s,x)\) is therefore comeager exactly when it is
\(X\), that is, exactly when \(\mathcal M\models\forall x\,\varphi\,[s]\).

Let \(\mu(\{a\})>0\) for every \(a\in X\). If \(N\) is \(\mu\)-null, choose
a Borel \(B\supseteq N\) with \(\mu(B)=0\). Any \(a\in B\) would give
\(\mu(B)\ge\mu(\{a\})>0\); so \(B=\varnothing\) and \(N=\varnothing\).
A section has \(\mu\)-measure one exactly when its complement is empty,
that is, when it is \(X\). The normalized counting measure gives each
point mass \(1/|X|>0\).

For the threshold binder, fix \(0<r<1\) and choose an integer
\(n\ge2\) with \(n\ge1/(1-r)\). Then \((n-1)/n\ge r\). Let \(R\) hold
at exactly \(n-1\) of \(n\) elements. The counting clause of
\cref{def:p15-pq-language} makes \((\exists^{\ge r}y)R(y)\) true, since
the proportion of witnesses is \((n-1)/n\), while \(\forall y\,R(y)\) is
false.
\end{proof}
\end{annexitem}

\begin{annexitem}{annex:p03-w03-p03-053-baire}{thm:p03-w03-p03-053-baire}{Nested balls}
\begin{proof}
Let \((X,d)\) be a nonempty complete metric space, let \(U\) be a nonempty
open set, and suppose \(U\subseteq\bigcup_{k<\omega}N_k\) with every
\(N_k\) nowhere dense. Put \(V_0=U\). Given a nonempty open \(V_k\), the
closed set \(\operatorname{cl}N_k\) has empty interior, so
\(V_k\not\subseteq\operatorname{cl}N_k\) and
\(V_k\setminus\operatorname{cl}N_k\) is a nonempty open set. Choose a
point \(x_{k+1}\) in it and a radius \(0<r_{k+1}<2^{-k}\) such that the
closed ball \(\bar B(x_{k+1},r_{k+1})\) lies in
\(V_k\setminus\operatorname{cl}N_k\), and put
\(V_{k+1}=B(x_{k+1},r_{k+1})\). For \(j>k\ge1\) we have
\(x_j\in V_k\), so \(d(x_j,x_k)<r_k<2^{1-k}\); the centres form a Cauchy
sequence and converge to some \(x\). For each \(k\ge1\) all \(x_j\) with
\(j\ge k\) lie in the closed ball \(\bar B(x_k,r_k)\), so \(x\) does too.
That ball lies in \(V_{k-1}\setminus\operatorname{cl}N_{k-1}\). Hence
\(x\in V_0=U\) and \(x\notin N_k\) for every \(k\), contrary to the
assumed covering. So \(U\) is not meager. (The choices are made by
dependent choice.)

A Polish space admits a complete metric. If \(\varnothing\) were comeager,
the nonempty open set \(X\) would be meager. A countable union of meager
sets is contained in a countable union of countably many nowhere dense
sets, which is again a countable union of nowhere dense sets; taking
complements, a countable intersection of comeager sets is comeager.
Finally let \(X\) be a nonempty countable Polish space. Every singleton is
closed. If no point were isolated, every singleton would have empty
interior and so be nowhere dense, and \(X=\bigcup_{x\in X}\{x\}\) would
be meager, which the first part excludes.
\end{proof}
\end{annexitem}

\begin{annexitem}{annex:p03-w03-p03-053-separation}{thm:p03-w03-p03-053-separation}{The set \(G\) and the two uniform laws}
\begin{proof}
\emph{Measure of a finitely determined set.} If \(A\subseteq2^\omega\) is
defined by prescribing the values at \(k\) given coordinates, all below
\(N\), then \(A\) is the disjoint union of the \(2^{N-k}\) basic sets
\([t]\) with \(|t|=N\) that agree with the prescription, so
\(\lambda(A)=2^{N-k}2^{-N}=2^{-k}\). Such a set is clopen, since its
complement is the union of the remaining basic sets of length \(N\).

Each \(B_n\) prescribes zeros at the \(n+1\) coordinates
\(n,\ldots,2n\), so it is clopen and \(\lambda(B_n)=2^{-(n+1)}\). The set
\(U_m=\bigcup_{n\ge m}B_n\) is open. It is dense: given a basic set
\([t]\), put \(n=\max(m,|t|)\) and extend \(t\) by zeros on all
coordinates from \(|t|\) to \(2n\) and arbitrarily afterwards; the result
lies in \([t]\cap B_n\). The complement of an open dense set is closed
with empty interior, hence nowhere dense. Therefore
\(2^\omega\setminus G=\bigcup_m(2^\omega\setminus U_m)\) is meager, and
\(G\) is comeager. As a countable intersection of open sets \(G\) is
Borel. Countable subadditivity gives
\(\lambda(G)\le\lambda(U_m)\le\sum_{n\ge m}2^{-(n+1)}=2^{-m}\) for every
\(m\), so \(\lambda(G)=0\).

In \((2^\omega,G)\) the section of \(G(x)\) is \(G\), which is comeager;
so \(Q_cx\,G(x)\) holds. The section of \(\neg G(x)\) is
\(2^\omega\setminus G\), whose complement \(G\) is a Borel set of measure
zero; so \(Q_mx\,\neg G(x)\) holds.

\emph{Uniform laws.} Write \(S_\varphi\) for the section of \(\varphi\)
at a fixed assignment; then \(S_{\varphi\land\psi}=S_\varphi\cap S_\psi\)
and \(S_{\neg\varphi}=X\setminus S_\varphi\). For \(Q_c\): if
\(S_\varphi,S_\psi\) are comeager, so is their intersection by
\cref{thm:p03-w03-p03-053-baire}; if \(S_\varphi\) and its complement were
both comeager, their empty intersection would be comeager, which that
lemma excludes. For \(Q_m\): if \(X\setminus S_\varphi\subseteq B_1\) and
\(X\setminus S_\psi\subseteq B_2\) with Borel null \(B_i\), then
\(X\setminus(S_\varphi\cap S_\psi)\subseteq B_1\cup B_2\), a Borel null
set; if \(S_\varphi\) and its complement both had measure one, both would
be null and \(X\) would be contained in a Borel null set, contradicting
\(\mu(X)=1\).

\emph{Mixed formula.} For \(\varphi=G(x)\) the antecedent
\(Q_cx\,G(x)\) is true and the consequent \(\neg Q_mx\,\neg G(x)\) is
false by the computation above.
\end{proof}
\end{annexitem}

\begin{annexitem}{annex:p03-w03-p03-053-invariance}{thm:p03-w03-p03-053-invariance}{Cardinal equality without measure or category equality}
\begin{proof}
\emph{An ordinary isomorphism.} The maps \(x\mapsto0^\frown x\) and
\(x\mapsto1^\frown x\) are bijections from \(2^\omega\) onto \(P_1\) and
onto \(2^\omega\setminus P_1\). The map sending \(x\) to the sequence
\(y\) with \(y(2k)=0\) and \(y(2k+1)=x(k)\) is a bijection from
\(2^\omega\) onto \(P_2\). The complement of \(P_2\) contains
\(\{y:y(0)=1\}\), which is in bijection with \(2^\omega\), and is a
subset of \(2^\omega\); by \cref{thm:p01-csb} it is in bijection with
\(2^\omega\). Composing these bijections gives a bijection
\(g:P_1\to P_2\) and a bijection
\(h:2^\omega\setminus P_1\to2^\omega\setminus P_2\); their union
\(f=g\cup h\) is a permutation of \(2^\omega\) with \(f[P_1]=P_2\), an
isomorphism of \(\{P\}\)-structures. An induction on
\(L_{\omega_1\omega}\)-formulas shows that an isomorphism preserves
satisfaction under transported assignments: atoms by the isomorphism,
negation and countable connectives directly, and quantifiers because
\(f\) is a bijection. Hence the two structures have the same
\(L_{\omega_1\omega}\)-theory.

\emph{Measure.} In the first structure the section of \(\neg P(x)\) is
the basic set \([1]\), whose complement \([0]\) has measure \(1/2\); every
Borel superset of \([0]\) has measure at least \(1/2\), so \([0]\) is not
null and \(Q_mx\,\neg P(x)\) is false. In the second, for each \(k\) the
set \(P_2\) is contained in the set prescribing zeros at the \(k\)
coordinates \(0,2,\ldots,2k-2\), which has measure \(2^{-k}\) by the
computation in \cref{annex:p03-w03-p03-053-separation}. The set \(P_2\)
is closed, as an intersection of clopen sets, so it is a Borel set of
measure zero and \(Q_mx\,\neg P(x)\) is true.

\emph{Category.} In the first structure the section of \(\neg P(x)\) is
comeager only if \([0]\) is meager, and \([0]\) is a nonempty open set,
which is not meager by \cref{thm:p03-w03-p03-053-baire}. In the second,
\(P_2\) is closed, and every basic set \([t]\) contains a sequence with
value \(1\) at some even coordinate \(2k\ge|t|\), which lies outside
\(P_2\). So \(P_2\) has empty interior, is nowhere dense, and its
complement, the section of \(\neg P(x)\), is comeager.

\emph{Preservation.} Let \(f:X\to Y\) be an \(L\)-isomorphism between
Polish-carrier structures. Prove by induction on formulas that
\(\mathcal M\models\varphi[s]\) iff \(\mathcal N\models\varphi[f\circ s]\)
for all assignments \(s\). Atoms, connectives and the ordinary quantifiers
are as in the first paragraph. For a new quantifier the induction
hypothesis gives \(f[\varphi(\mathcal M,s,x)]=\varphi(\mathcal N,f\circ s,x)\),
and \(f\) maps complements to complements because it is a bijection.
If \(f\) is a homeomorphism, it maps closures to closures and interiors to
interiors, hence nowhere dense sets to nowhere dense sets and meager sets
to meager sets; the same holds for \(f^{-1}\). So a section is comeager
iff its image is. If \(f\) and \(f^{-1}\) are Borel and
\(\nu(B)=\mu(f^{-1}[B])\) for Borel \(B\subseteq Y\), let \(A\subseteq C\)
with \(C\) Borel and \(\mu(C)=0\). Then \(f[C]=(f^{-1})^{-1}[C]\) is
Borel, and \(\nu(f[C])=\mu(C)=0\), so \(f[A]\) is \(\nu\)-null. For Borel
\(D\subseteq X\) we likewise have \(\mu(D)=\nu(f[D])\), so the same
argument applies to \(f^{-1}\). Thus a section has measure one iff its
image has measure one.
\end{proof}
\end{annexitem}

\begin{annexitem}{annex:p03-w03-p03-053-no-countable}{thm:p03-w03-p03-053-no-countable}{Isolated points and atoms}
\begin{proof}
At a value \(b\) for \(y\), the section of \(x\neq y\) is
\(X\setminus\{b\}\).

Let \(X\) be a countable Polish carrier. By
\cref{thm:p03-w03-p03-053-baire} it has an isolated point \(b\). Then
\(\{b\}\) is a nonempty open set, not meager by the same lemma, so
\(X\setminus\{b\}\) is not comeager and \(\sigma_c\) fails. Now let
\(\mu\) be a Borel probability measure on \(X\). Singletons are closed,
hence Borel, and countable additivity gives
\(1=\mu(X)=\sum_{x\in X}\mu(\{x\})\), so \(\mu(\{b\})>0\) for some \(b\).
Every Borel superset of \(\{b\}\) has measure at least \(\mu(\{b\})\),
so \(\{b\}\) is not null, \(X\setminus\{b\}\) does not have measure one,
and \(\sigma_m\) fails. Hence every model of either sentence has an
uncountable carrier.

In \(2^\omega\), each singleton \(\{b\}\) is closed, and it has empty
interior because every basic set contains at least two points; so it is
nowhere dense and \(2^\omega\setminus\{b\}\) is comeager. Also
\(\{b\}\subseteq[b\restriction k]\), so
\(\lambda(\{b\})\le2^{-k}\) for every \(k\), and \(\{b\}\) is a Borel null
set. Both sentences hold.

By \cref{thm:downward-ls}, applied with \(A=\varnothing\) in a
countable signature, every first-order sentence with a model has a model
of cardinality at most \(\aleph_0\). The sentences \(\sigma_c\) and
\(\sigma_m\) have models in regime (R5) and none of countable
cardinality, so that conclusion does not transfer.
\end{proof}
\end{annexitem}

\begin{annexitem}{annex:p03-w03-p03-053-l2-noncompact}{thm:p03-w03-p03-053-l2-noncompact}{Finiteness in the discrete topology}
\begin{proof}
In a topological structure satisfying \(\mathrm{Discrete}\), every point
\(x\) has an open \(U\) with \(x\in U\subseteq\{x\}\), so every singleton
is open and, taking unions, every subset is open. Open quantifiers of
\(L_2\) then range over all subsets of the carrier. With \(<\) a strict
linear order, the first conjunct inside \(\mathrm{Fin}\) says that a
nonempty \(U\) has a greatest element and the second that it has a least
element. So \(\mathrm{Fin}\) says that every nonempty subset has a greatest
and a least element.

\emph{Finite satisfiability.} A finite subset of \(\Gamma\) is contained,
for some \(N\), in
\[
 \{\mathrm{Lin},\mathrm{Discrete},\mathrm{Fin},
   \exists_{\ge1},\ldots,\exists_{\ge N}\}.
\] The set
\(\{0,\ldots,N-1\}\) with its usual order and the discrete topology
satisfies these: every nonempty subset is finite and so has a greatest
and a least element, by induction on its size.

\emph{No model.} Suppose a topological structure satisfies all of
\(\Gamma\). Its carrier \(X\) is infinite. Define
\(a_0=\min X\) and \(a_{k+1}=\min(X\setminus\{a_0,\ldots,a_k\})\); each
set is nonempty because \(X\) is infinite, and its least element is unique.
The set \(A=\{a_k:k<\omega\}\) is open, so by \(\mathrm{Fin}\) it has a
greatest element \(a_j\). But \(a_{j+1}\in A\), \(a_{j+1}\neq a_j\), and
\(a_j\) is the least element of a set containing \(a_{j+1}\), so
\(a_j<a_{j+1}\), a contradiction.

\emph{Proof systems.} Suppose \(\vdash\) were sound, strongly complete and
of finite character. Since \(\Gamma\) has no topological model, every
\(L_2\)-sentence, in particular \(\bot\), is a consequence of \(\Gamma\),
so \(\Gamma\vdash\bot\). Finite character gives a finite
\(\Gamma_0\subseteq\Gamma\) with \(\Gamma_0\vdash\bot\), and soundness
gives that \(\Gamma_0\) has no topological model, contrary to finite
satisfiability.
\end{proof}
\end{annexitem}

\begin{annexitem}{annex:p03-w03-p03-053-lt-reduction}{thm:p03-w03-p03-053-lt-reduction}{Basis axioms and \(L_t\)}
\begin{proof}
\emph{From two-sorted consequence to topological consequence.} Let
\(\mathcal T=(X,\tau,\ldots)\) be a topological model of \(\Gamma\). Read
it as a two-sorted structure with point sort \(X\), open sort \(\tau\)
and actual membership. Ext holds because open-sort identity is identity
of sets; Union because \(X\in\tau\); Empty because
\(\varnothing\in\tau\); Bas because \(U\cap V\in\tau\) is a witness.
By \cref{def:p03-topological-languages} the satisfaction clauses of
\(\mathcal T\) are the two-sorted first-order clauses of this structure,
so it satisfies \(\Gamma\cup\mathrm{Ax}\), hence \(\varphi\), and
\(\mathcal T\models\varphi\).

\emph{From topological consequence to two-sorted consequence.} Let
\(\mathcal S\) be a two-sorted model of \(\Gamma\cup\mathrm{Ax}\). Send
each open-sort element \(u\) to its extent \(\{x:x\varepsilon u\}\).
Ext makes this injective and reflects open equality; Union, Empty and Bas
make the extents a basis containing \(\varnothing\) that covers the
point sort and has the refinement property. With the point structure
unchanged, the extent map is an isomorphism of \(\mathcal S\) onto a
basoid model \(\mathcal B\); this is the construction in
\cref{annex:p03-topological-compactness-ls}. Isomorphic two-sorted
structures satisfy the same two-sorted sentences, so
\(\mathcal B\models\Gamma\). By \cref{thm:p03-topological-basis} the
generated topological model \(\widehat{\mathcal B}\) agrees with
\(\mathcal B\) on all \(L_t\)-sentences; as stated in
\cref{sec:p03-topological-reduced-products}, that proof is unchanged for
an arbitrary point signature. So \(\widehat{\mathcal B}\models\Gamma\), the
hypothesis gives \(\widehat{\mathcal B}\models\varphi\), and then
\(\mathcal B\models\varphi\) and \(\mathcal S\models\varphi\).

\emph{One-sorted encoding.} Add unary predicates \(\mathrm{Pt}\) and
\(\mathrm{Op}\). Translate a two-sorted formula \(\psi\) to \(\psi^*\)
by relativizing point quantifiers to \(\mathrm{Pt}\) and open
quantifiers to \(\mathrm{Op}\), and let \(\Delta\) say that both
predicates are nonempty, that they partition the domain, that every
constant lies in \(\mathrm{Pt}\), and that every point function maps
\(\mathrm{Pt}\)-arguments into \(\mathrm{Pt}\). A one-sorted model of
\(\Delta\) restricts to a two-sorted structure satisfying \(\psi\) iff it
satisfies \(\psi^*\), by induction on \(\psi\). Conversely a two-sorted
structure becomes a model of \(\Delta\) on the disjoint union of its
sorts, extending each point function to arguments outside \(\mathrm{Pt}\)
by a fixed point value and leaving relations unchanged; its restriction
is the original structure. Hence \(\Gamma\cup\mathrm{Ax}\) has
\(\varphi\) as a two-sorted consequence iff
\((\Gamma\cup\mathrm{Ax})^*\cup\Delta\) has \(\varphi^*\) as an
ordinary first-order consequence.

\emph{Consequences.} If \(\varphi\) is true in every topological model of
\(\Gamma\), the two previous paragraphs and
\cref{thm:fol-compactness} give a finite \(\Gamma_0\subseteq\Gamma\)
with \((\Gamma_0\cup\mathrm{Ax})^*\cup\Delta\models\varphi^*\), so
\(\varphi\) is true in every topological model of \(\Gamma_0\).
By \cref{thm:fol-completeness} the ordinary consequence is derivable
in \(\mathsf{ND}_{\mathsf{FOL}}\), which is sound; so this calculus,
applied to the encoded premises, is sound and strongly complete for
\(L_t\)-consequence over topological structures. With \(\Gamma\) empty,
\(\varphi\) is true in all topological structures iff
\(\mathrm{Ax}^*\cup\Delta\vdash\varphi^*\). For a recursive signature the
map \(\varphi\mapsto\varphi^*\) is computable, \(\mathrm{Ax}^*\cup\Delta\)
is finite apart from the recursive list of closure axioms for function
symbols, and derivations can be effectively enumerated; so the set of
such \(\varphi\) is recursively enumerable.
\end{proof}
\end{annexitem}

\begin{annexitem}{annex:p03-w03-p03-053-l2-uncountable}{thm:p03-w03-p03-053-l2-uncountable}{A complete order topology in \(L_2\)}
\begin{proof}
\emph{The real line is a model.} The usual order of \(\mathbb R\) is a
strict linear order, \(\mathbb R\) is nonempty, \((x+y)/2\) lies strictly
between \(x<y\), and \(x-1<x<x+1\); so \(\mathrm{Lin}\), \(\mathrm{Ne}\),
\(\mathrm{Dense}\) and \(\mathrm{NoEnd}\) hold. For \(\mathrm{Int}\), let
\(a,b\in\mathbb R\) and \(U=\{y:a<y<b\}\). If \(a\ge b\) then \(U\) is
empty, hence open. Otherwise every \(y\in U\) lies with the open interval
of radius \(\min(y-a,b-y)\) about it inside \(U\), so \(U\) is open in the
usual topology; in both cases \(U\) witnesses the open quantifier. For
\(\mathrm{Sup}\), let \(U\) be open, nonempty and bounded above by some
\(b\). By the least-upper-bound property of the real numbers \(U\) has a
supremum \(s\). Every \(y\in U\) satisfies \(y\le s\), that is
\(\neg\,s<y\); and if \(t<s\) then \(t\) is not an upper bound of \(U\),
so some \(y\in U\) has \(t<y\). Hence \(\mathrm{Sup}\) holds.

\emph{Every model is uncountable.} Let \((X,\tau,<)\) be a topological
model of \(\theta_{\mathbb R}\), and suppose towards a contradiction that
\(X\) is countable. By \(\mathrm{Ne}\) it is nonempty, so there is a map
\(k\mapsto q_k\) from \(\omega\) onto \(X\).

We define \(a_k,b_k\in X\) with \(a_k<b_k\) by recursion. By
\(\mathrm{Ne}\) choose \(x\in X\), by \(\mathrm{NoEnd}\) choose \(y<x\),
and put \(a_0=y\), \(b_0=x\). Given \(a_k<b_k\), apply \(\mathrm{Dense}\)
four times to obtain \(c_1,\ldots,c_4\) with
\(a_k<c_1<c_2<c_3<c_4<b_k\): \(c_1\) between \(a_k\) and \(b_k\), then
each \(c_{i+1}\) between \(c_i\) and \(b_k\). If \(q_k<c_2\) or
\(q_k=c_2\), put \((a_{k+1},b_{k+1})=(c_3,c_4)\); then
\(q_k<a_{k+1}\). Otherwise \(c_2<q_k\) by linearity; put
\((a_{k+1},b_{k+1})=(c_1,c_2)\); then \(b_{k+1}<q_k\). In both cases
\[
 a_k<a_{k+1}<b_{k+1}<b_k
 \qquad\text{and}\qquad
 q_k<a_{k+1}\ \text{or}\ b_{k+1}<q_k. \tag{\(\ast\)}
\]
For all \(j,m\) we have \(a_j<b_m\): if \(j\le m\) then
\(a_j\le a_m<b_m\), and if \(j>m\) then \(a_j<b_j\le b_m\), by the
monotonicity in \((\ast)\) and transitivity.

By \(\mathrm{Int}\), for each \(k\) there is \(U_k\in\tau\) whose members
are exactly the \(y\) with \(a_0<y<a_k\). Since \(\tau\) is a topology,
\(U=\bigcup_kU_k\in\tau\); its members are the \(y\) with \(a_0<y<a_k\)
for some \(k\). By \(\mathrm{Dense}\) some \(y\) has \(a_0<y<a_1\), so
\(U\) is nonempty; every member of \(U\) is below some \(a_k\), hence
below \(b_0\). So the hypothesis of \(\mathrm{Sup}\) holds for \(U\), and
there is \(s\in X\) such that no member of \(U\) lies above \(s\) and
every \(t<s\) lies below some member of \(U\).

First, \(a_j<s\) for every \(j\). By \(\mathrm{Dense}\) choose \(y\) with
\(a_j<y<a_{j+1}\). Then \(a_0\le a_j<y\), so \(y\in U_{j+1}\subseteq U\),
hence \(\neg\,s<y\), so \(y<s\) or \(y=s\) by linearity; in either case
\(a_j<s\).

Second, \(s<b_j\) for every \(j\). Otherwise \(b_j<s\) or \(b_j=s\), and
then \(b_{j+1}<b_j\) gives \(b_{j+1}<s\). So some \(y\in U\) has
\(b_{j+1}<y\). But \(y<a_i\) for some \(i\), and \(a_i<b_{j+1}\), so
\(y<b_{j+1}\), contradicting irreflexivity and transitivity of \(<\).

Hence \(a_{k+1}<s<b_{k+1}\) for every \(k\). By \((\ast)\),
\(s\neq q_k\) for every \(k\), so \(s\) is not in the range of the map
onto \(X\), a contradiction. Therefore \(X\) is uncountable.

\emph{Failure of the countable form.} The single sentence
\(\theta_{\mathbb R}\) has a topological model and no topological model
with a countable carrier. A countable-carrier form of downward
L\"owenheim--Skolem for \(L_2\) would supply one.
\end{proof}
\end{annexitem}

\begin{annexitem}{annex:p03-w03-p03-053-pure-equality}{thm:p03-w03-p03-053-pure-equality}{Equality types and quantifier rank}
\begin{proof}
All formulas in this proof are first-order formulas of the empty
signature, whose atoms are equations \(x=y\) between variables. Renaming
bound variables changes neither the quantifier rank nor the truth value
under any assignment, and \(\forall y\) is \(\neg\exists y\neg\); so we
may assume that each formula is built from atoms by \(\neg\), \(\land\)
and \(\exists\), and that a quantifier \(\exists y\) never binds a
variable free in the formula it is part of. Two tuples
\((a_1,\ldots,a_k)\) in \(A\) and \((b_1,\ldots,b_k)\) in \(B\) have the
\emph{same equality type} if \(a_i=a_j\iff b_i=b_j\) for all \(i,j\).

We prove by induction on \(\varphi\): \emph{if the free variables of
\(\varphi\) are among \(x_1,\ldots,x_k\), the quantifier rank of
\(\varphi\) is at most \(r\), \(|A|\ge k+r\), \(|B|\ge k+r\), and
\(\bar a\in A^k\), \(\bar b\in B^k\) have the same equality type, then
\(A\models\varphi[\bar a]\iff B\models\varphi[\bar b]\).}

For an atom \(x_i=x_j\) this is the definition of equality type.
Negation and conjunction do not raise the rank, so the claim for their
components with the same \(k\) and \(r\) gives it for the compound. Let
\(\varphi=\exists y\,\psi\), with the free variables of \(\psi\) among
\(x_1,\ldots,x_k,y\) and \(y\) distinct from these. Then \(\psi\) has
rank at most \(r-1\) and \(r\ge1\). Suppose
\(A\models\psi[\bar a,a]\) for some \(a\in A\). If \(a=a_i\) for some
\(i\), put \(b=b_i\). Otherwise choose \(b\in B\) distinct from
\(b_1,\ldots,b_k\); this is possible because \(|B|\ge k+r\ge k+1\). In
either case \((\bar a,a)\) and \((\bar b,b)\) have the same equality
type, and \(|A|,|B|\ge(k+1)+(r-1)\). The induction hypothesis for
\(\psi\), with \(k+1\) and \(r-1\), gives \(B\models\psi[\bar b,b]\), so
\(B\models\varphi[\bar b]\). The converse direction is symmetric.

For a sentence of rank at most \(q\) take \(k=0\) and \(r=q\); the empty
tuples have the same equality type, and the hypothesis is
\(|A|,|B|\ge q\).
\end{proof}
\end{annexitem}

\begin{annexitem}{annex:p03-w03-p03-053-finite-row}{thm:p03-w03-p03-053-finite-row}{Compactness, enumerability and interpolation over finite structures}
\begin{proof}
\emph{(1) Compactness.} Let \(\Gamma\) consist of the axioms of a strict
linear order for \(<\) and, for each \(m\ge1\), the first-order sentence
\(\exists_{\ge m}\) saying that there are at least \(m\) elements, as in
\cref{def:p03-finite-compactness-failure}. A finite subset of \(\Gamma\)
contains \(\exists_{\ge m}\) only for \(m\le M\) for some \(M\ge1\); the
usual order on \(\{1,\ldots,M\}\) is a nonempty finite model of it. A
nonempty finite structure with \(n\) elements fails
\(\exists_{\ge n+1}\). Hence \(\Gamma\models_{\mathrm{fin}}\bot\) while
no finite subset \(\Gamma_0\) has \(\Gamma_0\models_{\mathrm{fin}}\bot\).
By \cref{def:p15-pq-language}, \(L_{\omega P}\) retains first-order
formulas with their first-order satisfaction clauses, so the same
\(\Gamma\) shows the failure for \(L_{\omega P}\).

\emph{(2) Enumerability.} For first-order sentences this is
\cref{thm:trakhtenbrot}. For \(L_{\omega P}\), the sentences with rational
thresholds form a computably presented set of finite strings. The
first-order sentences form a decidable subset of it, those with no
probability binder, and on them the two satisfaction relations agree. So
the first-order sentences true in all nonempty finite structures are the
intersection of the \(L_{\omega P}\)-sentences with rational thresholds
true in all nonempty finite structures with a decidable set. If the latter
set were computably enumerable, so would the intersection be,
contradicting \cref{thm:trakhtenbrot}.

\emph{(3) Interpolation.} First we determine the finite models of
\(\varphi_1\). Let \(<\) be a strict linear order on a nonempty finite
set, listed as \(e_1<\cdots<e_n\). Then \(\mathrm{succ}(e_i,e_j)\) holds
exactly when \(j=i+1\): if \(j=i+1\) nothing lies strictly between; if
\(j>i+1\) the element \(e_{i+1}\) does; if \(j\le i\) then
\(e_i<e_j\) fails. The unique element with no predecessor is \(e_1\), and
the unique element with nothing above it is \(e_n\). Hence, if
\(\varphi_1\) holds, \(P(e_1)\), and by induction on \(i\) using the
alternation conjunct, \(P(e_i)\) holds exactly for odd \(i\); the last
conjunct gives \(P(e_n)\), so \(n\) is odd. Conversely, if \(n\) is odd
and \(P\) is interpreted as \(\{e_i:i\text{ odd}\}\), all four conjuncts
hold. The same argument for \(\varphi_2\) shows that a nonempty finite
structure satisfying \(\varphi_2\) has even size, and that on every set
of even size \(n\ge2\), ordered as \(e_1<'\cdots<'e_n\) with
\(P'=\{e_i:i\text{ odd}\}\), the sentence \(\varphi_2\) holds, because
then \(e_n\notin P'\).

A nonempty finite structure satisfying both \(\varphi_1\) and
\(\varphi_2\) would have a universe of odd size and of even size. So
\(\varphi_1\models_{\mathrm{fin}}\neg\varphi_2\).

Suppose \(\chi\) is a first-order sentence of the empty signature with
\(\varphi_1\models_{\mathrm{fin}}\chi\) and
\(\chi\models_{\mathrm{fin}}\neg\varphi_2\), and let \(q\) be its
quantifier rank. Put \(m=2q+1\). Let \(\mathcal A\) have universe
\(\{1,\ldots,m\}\), with \(<\) the usual order, \(P\) the odd numbers,
\(<'\) the usual order and \(P'=\varnothing\). Then
\(\mathcal A\models\varphi_1\), so \(\mathcal A\models\chi\). Let
\(\mathcal B\) have universe \(\{1,\ldots,m+1\}\), with \(<'\) the usual
order, \(P'\) the odd numbers, \(<\) the usual order and
\(P=\varnothing\). Then \(\mathcal B\models\varphi_2\). The truth of
\(\chi\) in a structure depends only on its universe, and
\(m,m+1\ge q\), so \(\mathcal B\models\chi\) by
\cref{thm:p03-w03-p03-053-pure-equality}. Then
\(\chi\models_{\mathrm{fin}}\neg\varphi_2\) gives
\(\mathcal B\models\neg\varphi_2\), a contradiction. So no such
\(\chi\) exists.
\end{proof}
\end{annexitem}

\begin{annexitem}{annex:p03-w03-p03-053-discrete-rank}{thm:p03-w03-p03-053-discrete-rank}{Cells, counts and quantifier rank on discrete spaces}
\begin{proof}
All formulas in this proof belong to \(L_2(\varnothing)\); their atoms are
\(x=y\), \(U=V\) and \(x\varepsilon U\). As in the proof of
\cref{thm:p03-w03-p03-053-pure-equality}, we may assume that formulas are
built from atoms by \(\neg\), \(\land\) and the two existential
quantifiers, and that no quantifier binds a variable free in the formula
it is part of; neither step changes quantifier rank or truth. All
structures are finite and discrete, so open quantifiers range over all
subsets of the carrier.

\emph{Configurations.} A configuration \((X;\bar a;\bar B)\) consists of a
finite set \(X\), a tuple \(\bar a=(a_1,\ldots,a_k)\in X^k\) and a tuple
\(\bar B=(B_1,\ldots,B_l)\) of subsets of \(X\), with \(k,l\ge0\). For
\(s\in\{0,1\}^l\) the \emph{cell} \(C_s(\bar B)\) is the set of
\(x\in X\) such that, for every \(j\le l\), \(x\in B_j\) exactly when
\(s(j)=1\). For \(l=0\) there is one cell, \(X\) itself. The cells are
pairwise disjoint and cover \(X\). The \emph{count} of a cell is
\(n_s=|C_s(\bar B)\setminus\{a_1,\ldots,a_k\}|\). Two configurations
\((X;\bar a;\bar B)\) and \((Y;\bar b;\bar D)\) of the same lengths
\(k,l\), with counts \(n_s\) and \(n'_s\), are \emph{\(r\)-similar} if
\begin{enumerate}
\item[(a)] \(a_i=a_j\iff b_i=b_j\) for all \(i,j\le k\);
\item[(b)] \(a_i\in B_j\iff b_i\in D_j\) for all \(i\le k\), \(j\le l\);
\item[(c)] \(\min(n_s,2^r)=\min(n'_s,2^r)\) for every \(s\).
\end{enumerate}
The relation is symmetric. If \(r'\le r\), then \(r\)-similar
configurations are \(r'\)-similar, since \(\min(n,2^{r'})=
\min(\min(n,2^r),2^{r'})\). If \(\min(n,2^r)=\min(n',2^r)\) and
\(n\neq n'\), then both minima equal \(2^r\), so \(n,n'\ge2^r\); and
\(n=0\iff n'=0\) because \(2^r\ge1\).

\emph{Splitting.} Let \(r\ge1\), \(h=2^{r-1}\), and let
\(\min(n,2h)=\min(n',2h)\) and \(n=d+e\) with natural numbers \(d,e\). Then
there are natural numbers \(d',e'\) with \(n'=d'+e'\),
\(\min(d,h)=\min(d',h)\) and \(\min(e,h)=\min(e',h)\). If \(n'=n\), take
\(d'=d\), \(e'=e\). Otherwise \(n,n'\ge2h\). If \(d<h\), take \(d'=d\) and
\(e'=n'-d\); then \(e=n-d>h\) and \(e'>h\), so both truncate to \(h\). If
\(e<h\), take \(e'=e\) and \(d'=n'-e\), symmetrically. If \(d,e\ge h\),
take \(d'=h\) and \(e'=n'-h\ge h\).

\emph{Claim.} If \((X;\bar a;\bar B)\) and \((Y;\bar b;\bar D)\) are
\(r\)-similar and \(\varphi\) has quantifier rank at most \(r\), free point
variables among \(x_1,\ldots,x_k\) and free open variables among
\(U_1,\ldots,U_l\), then \(X\models\varphi[\bar a,\bar B]\) if and only if
\(Y\models\varphi[\bar b,\bar D]\). We argue by induction on \(\varphi\).

\emph{Atoms.} For \(x_i=x_j\) use (a); for \(x_i\varepsilon U_j\) use (b).
For \(U_i=U_j\): \(B_i=B_j\) holds exactly when every cell \(C_s(\bar B)\)
with \(s(i)\ne s(j)\) is empty. A cell is empty exactly when its count is
\(0\) and no \(a_t\) lies in it. By (c), \(n_s=0\iff n'_s=0\). By (b),
\(a_t\in C_s(\bar B)\iff b_t\in C_s(\bar D)\), since both say that
membership in the \(j\)-th set agrees with \(s(j)\) for every \(j\). Hence
\(C_s(\bar B)=\varnothing\iff C_s(\bar D)=\varnothing\), and
\(B_i=B_j\iff D_i=D_j\).

\emph{Connectives.} \(\neg\) and \(\land\) do not raise the rank; apply the
hypothesis to the components with the same configurations.

\emph{Point quantifier.} Let \(\varphi=\exists x_{k+1}\psi\). Then
\(r\ge1\) and \(\psi\) has rank at most \(r-1\). Suppose
\(X\models\psi[\bar a,a,\bar B]\). If \(a=a_i\) for some \(i\), put
\(b=b_i\): (a) and (b) extend because \(a\) and \(b\) repeat \(a_i\) and
\(b_i\), and the counts are unchanged. Otherwise \(a\) lies in exactly one
cell \(C_s(\bar B)\) and is not among \(\bar a\), so \(n_s\ge1\), hence
\(n'_s\ge1\), and we choose \(b\in C_s(\bar D)\setminus\{b_1,\ldots,b_k\}\).
Then (a) extends because \(a\) and \(b\) are new, and (b) extends because
\(a\) and \(b\) lie in cells with the same index \(s\). The count of cell
\(s\) drops by one on each side and the others are unchanged. If
\(n_s=n'_s\), the new counts are equal; otherwise
\(n_s,n'_s\ge2^r\), so \(n_s-1,n'_s-1\ge2^r-1\ge2^{r-1}\). In both cases
the extended configurations are \((r-1)\)-similar, and the induction
hypothesis for \(\psi\) gives \(Y\models\psi[\bar b,b,\bar D]\). The
converse direction follows by symmetry.

\emph{Open quantifier.} Let \(\varphi=\exists U_{l+1}\psi\). Then \(r\ge1\)
and \(\psi\) has rank at most \(r-1\). Suppose
\(X\models\psi[\bar a,\bar B,E]\) for some \(E\subseteq X\). For each
\(s\) put \(d_s=|(C_s(\bar B)\cap E)\setminus\{a_1,\ldots,a_k\}|\) and
\(e_s=|(C_s(\bar B)\setminus E)\setminus\{a_1,\ldots,a_k\}|\), so
\(n_s=d_s+e_s\). By the splitting step with \(h=2^{r-1}\), choose
\(d'_s,e'_s\) with \(n'_s=d'_s+e'_s\), \(\min(d_s,h)=\min(d'_s,h)\) and
\(\min(e_s,h)=\min(e'_s,h)\). Choose
\(F_s\subseteq C_s(\bar D)\setminus\{b_1,\ldots,b_k\}\) with
\(|F_s|=d'_s\), and put
\[
 E'=\bigcup_s F_s\cup\{b_i: a_i\in E\}\subseteq Y .
\]
\(E'\) is open because \(Y\) is discrete. The cells of
\((\bar B,E)\) are the sets \(C_s(\bar B)\cap E\) and
\(C_s(\bar B)\setminus E\), with counts \(d_s\) and \(e_s\). On the other
side, \(F_s\subseteq C_s(\bar D)\), the cells \(C_s(\bar D)\) are pairwise
disjoint and no \(F_s\) meets \(\{b_1,\ldots,b_k\}\), so the cell
\(C_s(\bar D)\cap E'\) has count \(|F_s|=d'_s\) and the cell
\(C_s(\bar D)\setminus E'\) has count \(n'_s-d'_s=e'_s\). So (c) holds
with \(r-1\). For (b) at the new index: \(b_i\in E'\) exactly when
\(b_i=b_t\) for some \(t\) with \(a_t\in E\), since \(b_i\) lies in no
\(F_s\); by (a) this happens exactly when \(a_i=a_t\) for such a \(t\),
that is, when \(a_i\in E\). Conditions (a) and (b) for the old indices are
unchanged. The extended configurations are \((r-1)\)-similar, so
\(Y\models\psi[\bar b,\bar D,E']\) and \(Y\models\varphi[\bar b,\bar D]\).
The converse follows by symmetry.

\emph{Sentences.} For a sentence of rank at most \(q\) take \(k=l=0\).
The only cell is the carrier, with counts \(|X|\) and \(|Y|\), and
\(\min(|X|,2^q)=2^q=\min(|Y|,2^q)\) by hypothesis. So the empty
configurations are \(q\)-similar and the claim applies.
\end{proof}
\end{annexitem}

\begin{annexitem}{annex:p03-w03-p03-053-l2-interpolation}{thm:p03-w03-p03-053-l2-interpolation}{Finite carriers forced without a finiteness hypothesis}
\begin{proof}
\emph{(1) Models of \(\psi_1\).} Let \(\mathcal T=(X,\tau,<,P)\) be a
topological structure satisfying \(\psi_1\); its carrier is nonempty by
\cref{def:p03-topological-carrier}. By \(\mathrm{Discrete}\), every point
\(x\) has an open \(U\) with \(x\in U\subseteq\{x\}\); so every singleton
is open, every subset is a union of singletons, and \(\tau=\mathcal P(X)\).
The sentence \(\varphi_1\) contains \(\mathrm{Lin}\), so \(<\) is a strict
linear order, and \(\mathrm{Fin}\) then says that every nonempty subset of
\(X\) has a greatest and a least element. Suppose \(X\) were infinite. Put
\(a_0=\min X\) and \(a_{j+1}=\min(X\setminus\{a_0,\ldots,a_j\})\); each of
these sets is nonempty because \(X\) is infinite. The set
\(A=\{a_j:j<\omega\}\) has a greatest element \(a_j\). Since \(a_j\) is the
least element of a set that contains \(a_{j+1}\) and \(a_{j+1}\ne a_j\),
\(a_j<a_{j+1}\), a contradiction. So \(X\) is finite and nonempty. The
reduct \((X,<,P)\) is then a nonempty finite structure satisfying the
first-order sentence \(\varphi_1\), and by the determination of the finite
models of \(\varphi_1\) in the proof of
\cref{thm:p03-w03-p03-053-finite-row}(3), \(|X|\) is odd and \(P\) is the
set of points in odd positions of \(<\).

Conversely, let \(X\) be finite of odd size \(n\), listed as
\(e_1,\ldots,e_n\). Give it the discrete topology, the order
\(e_1<\cdots<e_n\) and \(P=\{e_i:i\text{ odd}\}\). Every point \(e_i\) has
the open set \(\{e_i\}\), so \(\mathrm{Discrete}\) holds. Every nonempty
subset of a finite linear order has a greatest and a least element, by
induction on its size, so \(\mathrm{Fin}\) holds. The first-order sentence
\(\varphi_1\) holds by \cref{thm:p03-w03-p03-053-finite-row}(3). The same
two arguments with \(<'\), \(P'\), \(\mathrm{Fin}'\) and \(\varphi_2\) give
the statement for \(\psi_2\), with even sizes \(n\ge2\) and
\(P'=\{e_i:i\text{ odd}\}\).

\emph{(2) Consequence.} The truth of an \(L_2(L^i)\)-sentence in a
topological structure depends only on the carrier, the topology and the
interpretations of the symbols of \(L^i\). So a topological
\((L^1\cup L^2)\)-structure satisfying \(\psi_1\land\psi_2\) has
\(L^1\)- and \(L^2\)-reducts satisfying \(\psi_1\) and \(\psi_2\), and by
(1) its carrier is finite of odd size and of even size. None exists, so
\(\psi_1\models_{\mathrm{top}}\neg\psi_2\).

\emph{(3) No interpolant.} Since \(L^1\cap L^2=\varnothing\), an
interpolant would be a sentence \(\chi\in L_2(\varnothing)\). Suppose
\(\psi_1\models_{\mathrm{top}}\chi\) and
\(\chi\models_{\mathrm{top}}\neg\psi_2\), and let \(q\) be the quantifier
rank of \(\chi\). Choose an odd \(m\ge2^q\). Let \(\mathcal A\) have carrier
\(\{1,\ldots,m\}\), the discrete topology, \(<\) the usual order, \(P\) the
odd numbers, \(<'\) the usual order and \(P'=\varnothing\). Its
\(L^1\)-reduct satisfies \(\psi_1\) by (1), so \(\mathcal A\models\psi_1\)
and \(\mathcal A\models\chi\). Let \(\mathcal B\) have carrier
\(\{1,\ldots,m+1\}\), the discrete topology, \(<'\) the usual order,
\(P'\) the odd numbers, \(<\) the usual order and \(P=\varnothing\). By
(1) for \(\psi_2\), \(\mathcal B\models\psi_2\). The truth of \(\chi\)
depends only on the carrier and the topology, and the two discrete spaces
have sizes \(m,m+1\ge2^q\), so \(\mathcal B\models\chi\) by
\cref{thm:p03-w03-p03-053-discrete-rank}. Then
\(\chi\models_{\mathrm{top}}\neg\psi_2\) gives
\(\mathcal B\models\neg\psi_2\), a contradiction. So no interpolant
exists. Since \(\psi_1\in L_2(L^1)\), \(\neg\psi_2\in L_2(L^2)\) and
\(\psi_1\models_{\mathrm{top}}\neg\psi_2\), the Craig interpolation
property of \cref{def:p03-w03-p03-053-l2-craig} fails.
\end{proof}
\end{annexitem}
